\section{Finite Fourier sums and Gauss's congruence}
\label{sec:fourier}

\paragraph{A weighted hyperbola identity.}
Fix $N\ge1$ and a table $w(1),\ldots,w(N)$ of integers. Partition the
pairs of positive integers $j,k$ with $jk\le N$ by their product. Each
pair in the part with product $u$ has the unique form $(u/k,k)$ with
$k\mid u$. Consequently
\[
   \sum_{jk\le N} w(jk)\mu(k)
    =\sum_{u=1}^N w(u)\sum_{k\mid u}\mu(k)=w(1).
\]
The constant table gives the floor-sum identity. The table
$w(u)=0,1,0,-1$ on residues $0,1,2,3\pmod4$ gives
$\sum_{jk\le N}\sin(\pi jk/2)\mu(k)=1$. These four integer values
supply the exact arithmetic interpretation of the sine. Each follows also
from the imaginary part of $i^u$.
Writing $W(k)=\sum_{1\le j\le N/k}w(jk)$ and $W(N+1)=0$, one more
finite rearrangement gives $\sum_{k=1}^N M(k)(W(k)-W(k+1))=w(1)$.
Indeed the coefficient of $\mu(d)$ after expansion telescopes to $W(d)$.
The same four-entry table therefore gives a Mertens version of the sine sum.


\paragraph{A polynomial certificate for the sum of primitive roots.}
Use rational polynomials stored as finite coefficient lists. For $n\ge1$
put $A_n=X^n-1$ and form the monic polynomials
\[
 D_n=\gcd\!\left(A_n,\prod_{e\mid n,\ e<n}(X^e-1)\right),
 \qquad H_n=A_n/D_n .
\]
An empty product is $1$. The gcd uses Euclidean division, which takes at
most one more than the degree of the smaller nonzero operand in division steps once the inputs are ordered by degree.
Since $A_n'=nX^{n-1}$ is coprime to $A_n$, the polynomial $A_n$ has no
nonconstant square divisor, since a repeated factor divides the derivative
by the product rule. For every proper divisor $e$ of $n$, the polynomial
$X^e-1$ divides both gcd operands and hence divides $D_n$. A common
factor of $H_n$ and $X^e-1$ would have its square dividing $H_nD_n=A_n$,
so their monic gcd is $1$. Here and below, one may verify coprimality by
the extended Euclidean algorithm. For $b>0$,
$\gcd(X^b-1,A_n)=X^{\gcd(b,n)}-1$ by repeated subtraction of
exponents, using
$X^a-1=X^{a-b}(X^b-1)+(X^{a-b}-1)$ for $a\ge b$.
In particular $\gcd(H_n,X^b-1)=1$ whenever $n\nmid b$.

Let $z$ denote the residue of $X$ modulo $H_n$, with equality tested by
polynomial division. Then $z^n=1$, and $z^b-1$ is invertible for every
positive $b$ with $n\nmid b$, by B\'ezout. For $h\ge0$, consider the
Ramanujan sum
\[
 C_n(h)=\sum_{0\le k<n,\ \gcd(k,n)=1}z^{hk}.
\]
Insert $[\gcd(k,n)=1]=\sum_{d\mid\gcd(k,n)}\mu(d)$. Writing $k=dj$ gives
an inner sum of $(z^{hd})^j$ for $0\le j<n/d$. If $n\mid hd$, each
summand equals $1$. Otherwise multiplication by the invertible $z^{hd}-1$
telescopes the sum to $z^{hn}-1=0$. Substitution of $e=n/d$ proves the
integer-valued formula
\[
       C_n(h)=\sum_{e\mid\gcd(n,h)}e\,\mu(n/e),
       \qquad C_n(1)=\mu(n).
\]
This is a uniform polynomial remainder identity, including $n=1$.
To evaluate at $\xi=\exp(2\pi i/n)$, each factor $\xi^e-1$ of the
proper-divisor product has modulus $2\sin(\pi\min(e,n-e)/n)>1/n$.
Indeed on $0\le t\le\pi/2<2$ the alternating sine series gives
$\sin t\ge t-t^3/6\ge t/3$. A product of $h$ such factors has modulus
at least $n^{-h}$. Writing the proper-divisor polynomial as $D_nV$,
the bound $|V(\xi)|\le 1+\sum_j|V_j|$ gives explicit apartness of
$D_n(\xi)$ from zero. Thus $A_n=D_nH_n$ forces $H_n(\xi)=0$.
Conversely every root of $H_n$ has exact order $n$ by the coprimalities
above. This evaluation gives the exponential formula, including the
convention $k=n$ in place of $k=0$.

\paragraph{The Farey formula.}
Let $\mathcal F_N$ consist of all reduced fractions in $[0,1]$ with
denominator at most $N\ge1$, both endpoints included. Sorting them has
no effect on a sum. Group the nonzero fractions by their denominator;
the group with denominator $d$ contains exactly the fractions $k/d$,
$1\le k\le d$, $\gcd(k,d)=1$. Its exponential sum is $\mu(d)$.
The extra term $0/1$ contributes $1$, yielding
\[
    \sum_{a\in\mathcal F_N}\exp(2\pi i a)=1+M(N).
\]
A single arithmetic certificate represents all summands in the residue
ring modulo $H_L$, $L=N!$. For denominator $d$, use $z^{L/d}$ in place
of the $d$th root. Its $b$th power equals $1$ if $d\mid b$, and
$z^{(L/d)b}-1$ is invertible otherwise. The same geometric-sum
calculation applies.

\paragraph{Generators modulo a prime.}
Let $p$ be prime and $m=p-1$. Multiplication by a nonzero residue $a$
permutes the $m$ nonzero residues of $\mathbb F_p$. Taking the product
and cancelling gives $a^m=1$. Its order is obtained by scanning
$1,\ldots,m$ for the first exponent giving $1$. Division with remainder
shows that this order divides every exponent giving $1$.
For each prime power $r^v$ exactly dividing $m$, the polynomial
$X^{m/r}-1$ has degree $m/r<m$; a nonzero field polynomial of degree $s$
has at most $s$ distinct roots, since division by $X-a$ removes a root
and reduces the degree by one. Scanning the nonzero residues therefore
returns an $a$ with $a^{m/r}\ne1$. The element $b_r=a^{m/r^v}$
has order $r^v$: its $r^v$th power is $1$, while its $r^{v-1}$st is
not. The product $g=\prod_r b_r$ has order $m$. Indeed if $g^u=1$,
raising to $m/r^v$ kills all the other factors and gives
$r^v\mid u(m/r^v)$, hence $r^v\mid u$. These divisibilities imply
$m\mid u$. For $m=1$ the empty product works.

The $m$ distinct powers $g^k$, $0\le k<m$, exhaust the nonzero
residues. Their orders are $m/\gcd(m,k)$, so the generators correspond
to the coprime exponents. Repeat the calculation of $C_m(1)$ with
$z=g$ in $\mathbb F_p$. Every nonzero difference is invertible by a
bounded field search. This proves Gauss's congruence
\[
  \sum_{\substack{1\le a<p\\ \operatorname{ord}_p(a)=p-1}} a
        \equiv \mu(p-1)\pmod p .
\]
In particular this yields $0$ when $p-1$ has a repeated prime factor,
and the parity sign otherwise.

\section{Counting irreducible polynomials over finite fields}
\label{sec:irreducibles}

Supply a finite field $F$ as a table of operations on $q\ge2$ labelled
elements, with distinguished $0,1$ and the field laws checked on the
table. A monic polynomial of degree $j$ is encoded by its $j$ free
coefficients, a base-$q$ word numbered from $0$ through $q^j-1$.
Let $I_q(j)$, $j\ge1$, count those with no monic divisor of degree $1$ through
$\lfloor j/2\rfloor$. The scan is exact since any proper factorization
has a smaller factor in this range and the leading coefficient can be
normalized. This defines a primitive recursive count.

\paragraph{Unique factorization suffices.}
A monic polynomial of positive degree either passes the irreducibility
test or splits into two monic factors of smaller positive degree. Repeating
the split gives irreducibles: the number of factors increases at each
split and cannot exceed the original degree. Over the field table,
long division cancels one leading term per step. Iterated division with
remainders, recording the two linear-combination coefficients, gives
B\'ezout's identity for a monic gcd. Thus if an irreducible $P$ divides
$AB$ and does not divide $A$, an identity $UP+VA=1$ multiplied by $B$
shows that $P$ divides $B$. It follows by stripping one irreducible at
a time and cancelling it that the multiset of monic irreducible factors
is unique. Every decision here is a comparison of finite codes.

\paragraph{The coefficient identity.}
Fix $n\ge1$ and work through degree $n$ in the ring
$\mathbb Z[T]/(T^{n+1})$. Enumerate all monic irreducibles $P$ of
degree $1$ through $n$ and put
\[
  B(T)=\sum_{j=0}^{n}q^jT^j
     =\prod_P (1-T^{\deg P})^{-1}.
\]
Each inverse is a finite geometric series in this ring. The coefficient
of $T^j$ in the product counts multisets of irreducibles of total degree
$j$, hence monic polynomials of degree $j$ by the preceding bijection.
The degree-zero monic polynomial is $1$. On this truncated ring the operator
$\Delta=T\,d/dT$ satisfies the product rule: it multiplies the $j$th
coefficient by $j$ and preserves the truncation. All factors in $B$ have
constant term one, so are units. Applying $\Delta B/B$ on both sides,
using $(1-qT)B=1$, gives
\[
  \frac{qT}{1-qT}
     =\sum_{d=1}^n I_q(d)\frac{dT^d}{1-T^d}.
\]
Comparison in degree $n$ gives $q^n=\sum_{d\mid n} d I_q(d)$.
Divisor inversion now yields
\[
             n I_q(n)=\sum_{d\mid n}\mu(d)q^{n/d}.
\]
The integer on the right is divisible by $n$ and nonnegative.
The proof uses only products of bounded length and coefficient arrays.

\paragraph{Constructing a field for every prime-power input.}
For completeness take $q=p$ prime first, with the residue field table
obtained by modular arithmetic. The coefficient count implies
$d I_p(d)\le p^d$ for $d\ge1$. For $r\ge2$ the contribution of all
proper divisors of $r$ is at most
$\sum_{d=1}^{\lfloor r/2\rfloor}p^d<p^r$; this follows by summing the
geometric progression and using $\lfloor r/2\rfloor+1\le r$.
Thus $I_p(r)>0$; for $r=1$ it equals $p$. A bounded scan of the $p^r$
monic degree-$r$ polynomials supplies an irreducible $P$. Remainders
of degree $<r$ modulo $P$ form the $p^r$-element field: B\'ezout
computes an inverse for every nonzero residue. The count just proved
can therefore be instantiated at each $q=p^r$ without asking for an
extension field or an irreducible polynomial as an oracle.

Conversely the size of every supplied finite field is a prime power. Among
the $q+1$ multiples $0,1,2\\cdot1,\\ldots,q\\cdot1$, two agree, so a
bounded scan gives a least positive $p$ with $p\\cdot1=0$. If $p=uv$
with $1<u,v<p$, the nonzero elements $u\\cdot1$ and $v\\cdot1$ have zero
product, contradicting the field laws. The multiples of $1$ form a copy
of $\\mathbb F_p$. Starting from the span $\\{0\\}$, scan the field for the
first element $z$ outside the current span $S$. Adjoining $z$ produces
the $p$ disjoint translates $S+az$ with $a\\in\\mathbb F_p$: equality
across two translates would express $z$ in $S$ by division by a nonzero
scalar. Each step multiplies the size by $p$, and after at most $q$ steps
the span is the whole field. Thus $q=p^r$ with $r\\ge1$.

\section{Incidence inversion and convolution powers}
\label{sec:incidence}
For a finite poset given by its decidable order table, fix an interval
$[a,b]$. Set $u(a,a)=1$ and, for $a<y\le b$, put
$u(a,y)=-\sum_{a\le x<y}u(a,x)$. To evaluate the recurrence,
order the elements by increasing cardinality of $[a,y]$, breaking ties
by their stored indices. Strict predecessors always have smaller
cardinality. Thus $\sum_{a\le x\le y}u(a,x)=[a=y]$, and each last
summand is determined by those preceding it. Equivalently, if the
order matrix in a linear extension is $Z=I+U$ on $m$ elements, the
strictly triangular matrix $U$ satisfies $U^m=0$. Multiplication shows
that $\sum_{j=0}^{m-1}(-U)^j$ is an inverse of $Z$ on both sides.
Its entries are the same $u$, with zero outside
comparable pairs. The proof applies to supplied finite intervals of a
locally finite poset whenever their order tables are given.
In the positive-integer divisibility order, the interval from $a$ to
$b$ with $a\mid b$ has exactly the elements $ad$ for $d\mid b/a$.
Hence $u(a,b)=\mu(b/a)$ by divisor cancellation and uniqueness.

For $k\ge0$, let $\mu_k$ be the $k$-fold Dirichlet convolution of
$\mu$, with $\mu_0(n)=[n=1]$. If $f,g$ are multiplicative then so is
$f*g$. Every divisor of a coprime product $ab$ corresponds uniquely to
a pair of divisors of $a,b$ by taking prime factors from the two disjoint
supports. The sum defining $(f*g)(ab)$ therefore separates into the two
sums at $a$ and $b$. At a prime power $p^v$, a nonzero term of the
$k$-fold convolution distributes the $v$ copies of $p$ among distinct
slots. Choosing the slots gives
\[
  \mu_k(p^v)=(-1)^v\binom{k}{v},\qquad
  \mu_k(n)=\prod_{p^v\parallel n}(-1)^v\binom{k}{v}.
\]
Here $p^v\parallel n$ in the product means $v>0$ is the exact exponent
and $\binom{k}{v}=0$ for $v>k$. For the polynomial extension, replace
$\binom{k}{v}$ by
$B_v(K)=K(K-1)\cdots(K-v+1)/v!$ in $\mathbb Q[K]$ with $B_0=1$.
The parameter identity $\mu_{K+L}=\mu_K*\mu_L$ follows locally from
$B_v(K+L)=\sum_{j=0}^v B_j(K)B_{v-j}(L)$: at nonnegative integer
arguments the equation counts $v$-element subsets of two disjoint
blocks. Each variable has degree at most $v$, so the finite grid
$0,\ldots,v$ in each variable proves equality of polynomials by the
root bound, applied first in one variable and then the other.
Prime-factor multiplicativity gives the general identity. These are
rational-coefficient array equalities in iPRA and specialize to complex
parameters under evaluation.
